\section{Preliminaries}



\subsection{Notation}

For integers $m\leq n$, let $[m,n]:=\{m,\ldots,n\}$ and
$[m,n):=\{m,\ldots,n-1\}$. In particular,
$[0,n)=\{0,1,\ldots,n-1\}$, whereas $[n]:=[1,n]$ is one-based.
We write index ranges explicitly when addressing bins, blocks, or tree nodes.
The computational security parameter is $\kappa$; $\lambda$ controls
statistical error. Our concrete examples use $\kappa=128$ and $\lambda=40$.
The parties are $P_0,P_1$.
We use $:=$ for deterministic assignment and $\gets$ for sampling.
An indicator $[E]$ is one when predicate $E$ holds and zero otherwise.



We denote a sharing of $x$ by $\share x$ and the share held by $P_i$
by $\share x_i$. Binary sharing means $x=\share x_0\oplus\share x_1$.
For a finite additive group $\mathbb G$, additive sharing means
$x=\share x_0+\share x_1\in\mathbb G$.
The notation $\share x\in\mathbb G$ specifies the reconstructed value's
domain. A shared public constant uses a fixed local sharing known to both
parties. Ideal circuit and permutation calls return fresh random shares
unless their interfaces state otherwise. The DPF and DMPF functionalities
instead let the corrupt party choose its output share and return the
complementary share to the honest party.

\subsection{Functions}
We make use of the following helper functions (inputs may be secret-shared unless noted).

\begin{itemize}
  \item $\decomp_w:[0,w^D)\to[0,w)^D$ returns the $D$ base-$w$ digits
  of an integer, for public integers $w\geq2$ and $D\geq1$.
  Thus $a=\sum_{i=1}^D w^{i-1}b_i$. Protocol boxes specify the digit order
  when assigning these outputs.

  \item $\unitVec_{S}:\ (\,\alpha\in S,\ \beta\in \GG\,)\ \rightarrow\ \GG^{S}$, where
  $v=\unitVec_{S}(\alpha,\beta)$ satisfies $v_\alpha=\beta$ and $v_j=0$ for $j\in S\setminus\{\alpha\}$.

  \item $\append:X^n\times X^m\to X^{n+m}$ concatenates two sequences.
  We also write $x\Vert y$ for the same operation.

  \item $\sort:\ X^{n}\times \{\text{total order }\prec\}\ \rightarrow\ (X^{n},\pi)$.
  Given a sequence and a comparator, outputs the \emph{stably} sorted sequence (ascending under $\prec$) and the permutation
  $\pi\in S_n$ such that $\text{sorted}=x_{\pi(1)},\ldots,x_{\pi(n)}$. For tuples, the default is lexicographic order on the first component.

  \item $\convert_\GG:\ \{0,1\}^{\ell}\ \rightarrow\ \GG$ is a fixed, deterministic map used to turn a leaf’s $\ell$-bit string into a group element.
  Two standard instantiations we rely on:
  (i) if $\GG=(\ZZ_p,+)$ or $\FF_p$, interpret the bit string as an integer and reduce modulo $p$;
  (ii) for an arbitrary prime-order group, use a fixed hash-to-group encoding $H:\{0,1\}^{\ell}\!\to\!\GG$ (e.g., a standard hash-to-field followed by a canonical map).
  The choice is public and fixed by parameters; security only requires that outputs are computationally indistinguishable from random in~$\GG$ given random inputs.

  \item $\shuffle:\ X^{n}\ \rightarrow\ X^{n}$ samples a uniform $\pi\gets S_n$ and returns
  $(x_{\pi(1)},\ldots,x_{\pi(n)})$.
  On shared inputs, this denotes an ideal functionality: reconstruct the
  records internally, sample $\pi$, and return fresh shares of the permuted
  records. Neither party receives $\pi$.
  A named reusable instance $\mathsf{Shuf}$ samples and retains $\pi$ on
  $\mathsf{Setup}$; each $\mathsf{Expand}$ applies that same permutation to a
  new shared vector and returns fresh shares. Both methods return the
  permuted vector; repeated setup or expansion before setup returns $\bot$.
  Aligned record fields use the same permutation.

  \item $\mathsf{reshare}(\share x\in\mathbb G)\to\share x$ denotes the
  ideal identity circuit with fresh additive output shares. It reconstructs
  $x$ internally, samples $r\gets\mathbb G$, and returns $r,x-r$.
  No reconstructed value is opened.

  \item $\mathsf{permute}$ denotes an ideal privately controlled permutation:
  one designated party supplies $\pi$, the parties supply shared records,
  and the functionality returns fresh shares of their permutation.
  Only the designated party knows $\pi$. A named instance retains it for
  inverse calls. Composing one such instance per party gives the two
  serial hidden passes used by Waterfall.

  \item $\zip:\ X_1^{n}\times \cdots \times X_k^{n}\ \rightarrow\ (X_1\times\cdots\times X_k)^{n}$, denote the ``zip'' operator
  $\zip(\vec{x}^{(1)},\ldots,\vec{x}^{(k)})=((x^{(1)}_1,\ldots,x^{(k)}_1),\ldots,(x^{(1)}_n,\ldots,x^{(k)}_n))$.
  The inverse $\uzip$ unzips a list of $k$-tuples back into $k$ aligned arrays of length $n$. We will also use them as infix binary operators. 

  \item $\powerset:\ 2^{U}\rightarrow 2^{2^{U}}$ is the usual power-set operator, returning the set of all subsets of~$U$.

  \item $\G:\{0,1\}^{\kappa}\to(\{0,1\}^{\kappa})^w$ is the tree PRG
  for branching factor $w$. Each child seed includes its tag as its least
  significant bit. The protocol extracts these bits; $\G$ has no separate
  auxiliary output.
\end{itemize}



\subsection{Related Work}

\paragraph{Function secret sharing and DPFs.}
Function Secret Sharing (FSS) provides short keys that secret-share the evaluation of a function across two (or more) parties. The special case for point functions—Distributed Point Functions (DPFs)—originates with Gilboa-Ishai and subsequent improvements \cite{dpf,dpf2}. Later works explored programmability and larger payloads, along with protocols for efficient \emph{distributed} key generation without a trusted dealer \cite{boyle2022programmable,dpfKeygen}. These lines supply the core building blocks we rely on (standard and sparse DPFs) and motivate our fully distributed DMPF keygen.

\paragraph{DPFs in PIR, anonymity, and aggregation.}
DPFs have been used as a backbone for two-server PIR and anonymous write/read systems (e.g., Riposte), where a client encodes a unit vector query/update via DPF keys \cite{corrigan2015riposte}. In private telemetry and measurement, Prio and successors (VDAF) use DPF-style ideas to aggregate client data while preserving privacy and enabling robustness/verification \cite{prio,vdaf}. Recent systems also study aggregation of high-dimensional sparse vectors with DPF-based encodings \cite{asi2025preamble} and incremental/offline-online PIR updates compatible with DPF pipelines \cite{inc-offline-online-pir}. Our setting differs: we compress and \emph{multiply} sparse structures arising inside PCGs—hence the need for DMPFs rather than per-point DPF summation.

\paragraph{DMPFs and big-state constructions.}
DMPFs can be viewed as an instance of earlier FSS-for-batch-codes constructions \cite{fssBatchCodes}. More recently, DMPFs were explicitly identified and optimized, with cuckoo-style schemes and big-state variants providing sublinear domains for sub-instances \cite{dmpf}.  Our Waterfall construction retains the cuckoo-style sparse layout but uses an input-independent candidate matrix with a fixed scan and bounded complete repair.  We also give bucketing and big-state key-generation paths, and later isolate Reverse Cuckoo as a Ring-LPN-specific optimization rather than a generic DMPF.

\paragraph{PCGs from LPN/Ring-LPN and beyond.}
PCGs unify many preprocessing correlations for MPC and were introduced as a unifying abstraction by Boyle et al.\ \cite{pcg2019}. Ring-LPN PCGs first achieved silent OLE/triples from structured noise \cite{boyle2020efficient}, with later generalizations to arbitrary fields and to $\mathbb{Z}/p^k\mathbb{Z}$ \cite{anyField,pcgZ2k}. Beyond LPN-based assumptions, Bombar et al.\ introduced programmable PCGs from the quasi-abelian syndrome decoding (QA-SD) problem \cite{pcgQASD}. Recent high-throughput Boolean PCGs (e.g., FOLEAGE) highlight how practical pipelines are gated by multi-point encodings in their sparsity steps \cite{foleage}. Our DMPFs reduce the number of domain traversals used to encode sparse vectors. Transferring this improvement to a particular PCG requires checking its noise distribution and security assumptions.

\paragraph{Implementation of PCGs.}
Silentium \cite{silentium} implements a PCG for Beaver triples.
Our evaluation instead measures the Reverse-Cuckoo Ring-LPN pipeline over
two prime fields, with the baselines and setup exclusions stated in
\sectionref{sec:eval}. We do not claim the same gains for every PCG.

\paragraph{Cuckoo hashing and stash bounds.}
We rely on classical cuckoo hashing analysis and its stash refinements to bound the correctness error of fixed Waterfall configurations \cite{cuckoo,kirsch2010more,DBLP:journals/corr/abs-1204-4431}. For set-based cryptographic protocols such as PSI, cuckoo hashing is now standard, and circuit/OPPRF-based PSI variants analyze related load and failure profiles \cite{pinkas2018efficient,demmler2018pir,yeo2025private}.
